\documentclass[12pt]{ai_conquer2026}
\addbibresource{references.bib}

\title{%
  \begin{center}
        \Large {AI VIET NAM -- AI Conquer 2026}\\[.5ex]  
        \Huge \textbf{Guidance to AI Conquer 2026 \\ 
        (\LaTeX\space Blog Template)}\\[.5ex]
  \end{center}
}

% Author --------------------------------------------------------------------
\posttitle{
    \author{
        \begin{minipage}{\textwidth}
        \centering
        {
        \vspace{-1.5em}
          Ban biên soạn AIO
        }\\
      \end{minipage}
    }
}
\date{}

% Header --------------------------------------------------------------------
\pagestyle{fancy}
\fancyhf{}
\lhead{\href{https://www.facebook.com/aivietnam.edu.vn}{\textcolor{aivnGreen}{\bfseries AI VIET NAM (AI Conquer 2026)}}}
\rhead{\href{https://aivietnam.edu.vn/}{\textcolor{aivnGreen}{\bfseries aivietnam.edu.vn}}}
\fancyfoot[C]{\thepage}
\renewcommand{\headrulewidth}{1.0pt}
\renewcommand{\headrule}{{\color{aivnGreen}\hrule width\headwidth height\headrulewidth \vskip-\headrulewidth}}

\begin{document}
\maketitle
\vspace{-5.5em}
\setlength{\epigraphwidth}{0.6\textwidth}

\section{Tổng quan bài viết}

Phần này dùng để giới thiệu ngắn gọn bài viết blog. Nội dung nên cho biết bài viết nói về chủ đề gì, dành cho ai, và người đọc sẽ thu được điều gì sau khi đọc xong. Đây là phần mở đầu giúp xác định giọng điệu và trọng tâm của bài viết trước khi đi vào vấn đề, cách tiếp cận, triển khai, và kết quả.

\newpage
\setcounter{tocdepth}{2}
\tableofcontents
\thispagestyle{fancy}

\newpage
\section{Giới thiệu}

Phần này mở đầu bài viết bằng cách giới thiệu chủ đề và tạo ngữ cảnh cho người đọc. Nội dung nên ngắn gọn nhưng đủ hấp dẫn để người đọc hiểu vì sao chủ đề này đáng quan tâm. Nếu phù hợp, có thể bắt đầu từ một tình huống thực tế, một câu hỏi, hoặc một vấn đề quen thuộc để dẫn vào nội dung chính.

\section{Vấn đề}

Phần này trình bày vấn đề mà bài viết muốn giải quyết hoặc phân tích. Người viết nên mô tả vấn đề một cách rõ ràng, nêu được khó khăn cốt lõi, và chỉ ra tác động của nó trong thực tế. Đây là phần giúp người đọc hiểu vì sao cần đến cách tiếp cận hoặc giải pháp được trình bày ở các phần sau.

\section{Cách tiếp cận}

Phần này mô tả ý tưởng chính hoặc hướng tiếp cận được lựa chọn để xử lý vấn đề. Nội dung nên tập trung vào logic tổng thể, lý do chọn cách làm đó, và bức tranh khái quát của giải pháp. Với blog, phần này nên giữ ngôn ngữ dễ tiếp cận, tránh đi quá sâu vào chi tiết kỹ thuật nếu không thật sự cần thiết.

\section{Triển khai}

Phần này trình bày cách ý tưởng được hiện thực hóa. Nội dung có thể bao gồm các bước xây dựng, công cụ được dùng, những quyết định quan trọng trong quá trình triển khai, hoặc các khó khăn phát sinh khi làm thực tế. Đây là phần giúp người đọc hiểu cách biến một hướng tiếp cận thành một giải pháp hoạt động được.

\section{Kết quả}

Phần này trình bày kết quả đạt được sau khi triển khai. Nội dung có thể là kết quả định lượng, đầu ra minh họa, so sánh trước và sau, hoặc những quan sát quan trọng rút ra từ quá trình thực hiện. Mục tiêu là giúp người đọc thấy được giá trị thực tế của giải pháp thay vì chỉ dừng ở mặt ý tưởng.

\section{Bài học rút ra}

Phần này tổng hợp những kinh nghiệm, nhận xét, hoặc bài học quan trọng sau quá trình thực hiện. Người viết có thể nêu điều gì đã làm tốt, điều gì chưa hiệu quả, và nếu làm lại thì nên thay đổi ở đâu. Đây là phần rất phù hợp với blog vì nó tạo cảm giác chân thực và giúp người đọc học được điều gì đó vượt ra ngoài kết quả cuối cùng.

\section{Kết luận}

Phần cuối dùng để khép lại bài viết. Nội dung nên tóm tắt ngắn gọn vấn đề, cách tiếp cận, và kết quả chính, đồng thời để lại một thông điệp rõ ràng cho người đọc. Nếu phù hợp, có thể thêm một gợi ý về hướng phát triển, câu hỏi mở, hoặc lời mời người đọc tiếp tục tìm hiểu sâu hơn về chủ đề.

\newpage
\section{Tài liệu tham khảo}
\label{section:references}
\nocite{*}
\vspace{-3.2em}
\printbibliography

\newpage
\begin{appendices}
\begin{enumerate}
    \item \textbf{AIVIETNAM:} Thông tin về AIVIETNAM có thể đọc thêm tại \href{https://aivietnam.edu.vn/}{đây}.
\end{enumerate}
\end{appendices}

\end{document}